% STATUS: rewritten 2026-11-05
% STATUS: extended 2026-11-07 (Abramov and Sovetnikov 2025 read at the
% introduction, Definition 2.1 and Theorem 4)
% Sources read in full: Li 2018 (definitions and divergence, p. 657; Theorem 2.1,
% pp. 658--659); Abramov and Sovetnikov 2025 (pp. 487--489, 495).
\section{Poisson algebras}\label{sec:nonassociative:poisson}

Every derivation below acts on a Poisson algebra, where the product rule must hold for the commutative product and the bracket together.

The $\delta$-derivations of the previous section relax the product rule by a
scalar. % claim-ok: own summary of the previous section
A Poisson algebra instead carries two products, and a derivation must
respect both \cite[p.~657]{AlbertiArestiLi2018Poisson}.

\begin{definition}[{\cite[p.~657]{AlbertiArestiLi2018Poisson}}]
\label{def:nonassociative:poisson}
A vector space $B$ over a field $k$ is a \emph{Poisson algebra} if $B$
carries two bilinear operations, a multiplication $x\cdot y$ and a bracket
$\{x,y\}$, such that $B$ is a commutative associative algebra under the
multiplication, $B$ is a Lie algebra under the bracket, and
\[
\{x\cdot y,z\}\;=\;x\cdot\{y,z\}+\{x,z\}\cdot y
\qquad\text{for all }x,y,z\in B.
\]
\end{definition}

The last identity says that the map $x\mapsto\{x,z\}$, for a fixed $z$,
is a derivation of the associative product. % claim-ok: own restatement of the Leibniz rule above
The parallel requirement defines the derivations of
the structure. % claim-ok: own linking sentence

\begin{definition}[{\cite[p.~657]{AlbertiArestiLi2018Poisson}}]
\label{def:nonassociative:poisson-derivation}
Let $B$ be a Poisson algebra. A linear map $D$ on $B$ is a \emph{Poisson
derivation} if
\[
D(xy)\;=\;D(x)y+xD(y)
\quad\text{and}\quad
D(\{x,y\})\;=\;\{D(x),y\}+\{x,D(y)\}
\qquad\text{for all }x,y\in B.
\]
\end{definition}

The basic example is built on a polynomial algebra. % claim-ok: own linking sentence

\begin{example}[{\cite[p.~657]{AlbertiArestiLi2018Poisson}}]
\label{ex:nonassociative:symplectic}
The \emph{symplectic Poisson algebra} $S_n$ is the polynomial algebra
$k[x_1,y_1,\dots,x_n,y_n]$ with the bracket given by
\[
\{x_i,y_j\}=\delta_{ij},\qquad \{x_i,x_j\}=0,\qquad \{y_i,y_j\}=0,
\]
where $\delta_{ij}$ is the Kronecker symbol and $1\le i,j\le n$.
\end{example}

A derivation $D$ of a polynomial algebra in $n$ generators can be written
as $D=\sum_{i=1}^{n}D(x_i)\,\partial/\partial x_i$
\cite[p.~657]{AlbertiArestiLi2018Poisson}. Li denotes by
\[
D^{\ast}\;=\;\sum_{i=1}^{n}\frac{\partial}{\partial x_i}D(x_i)
\]
the \emph{divergence} of $D$ \cite[p.~658]{AlbertiArestiLi2018Poisson}.
% notion-ok: divergence as named at \cite[p.~658]{AlbertiArestiLi2018Poisson}
For a
Poisson derivation of the symplectic algebra the divergence vanishes, and
more holds (Theorem~\ref{thm:nonassociative:poisson-inner}).

\begin{theorem}[{\cite[Theorem~2.1, pp.~658--659]{AlbertiArestiLi2018Poisson}}]
\label{thm:nonassociative:poisson-inner}
If $D$ is a derivation of the symplectic Poisson algebra
$S_n=k[x_1,y_1,\dots,x_n,y_n]$, then $D^{\ast}=0$ and there exists
$f\in S_n$ such that $D=\operatorname{ad}_f$, that is, $D(g)=\{f,g\}$ for
any $g\in S_n$.
\end{theorem}

Every Poisson derivation of $S_n$ is therefore inner. A transposed Poisson
algebra trades the Leibniz rule of the bracket for a factor one half. Let $P$
carry a commutative associative product and a Lie bracket such that
$2x\cdot[y,z]=[x\cdot y,z]+[y,x\cdot z]$ for all $x,y,z$
\cite[equation~(12), p.~489]{AbramovSovetnikov2025}. A transposed Poisson
superalgebra carries the graded form of the same identity
\cite[Definition~2.1, p.~490]{AbramovSovetnikov2025}. Abramov and Sovetnikov use
an even derivation to build new structures from it (Theorem~\ref{thm:nonassociative:transposed-poisson-3lie})
\cite[p.~487]{AbramovSovetnikov2025}.\begin{theorem}[{\cite[Theorem~4, p.~495]{AbramovSovetnikov2025}}]\label{thm:nonassociative:transposed-poisson-3lie}
Let $(P,\cdot,[\,,\,])$ be a transposed Poisson superalgebra and let $D$ be its
even derivation. Define
$[x,y,z]=D(x)\cdot[y,z]+(-1)^{|x||yz|}D(y)\cdot[z,x]+(-1)^{|xy||z|}D(z)\cdot[x,y]$.
Then $(P,[\,,\,,\,])$ is a 3-Lie superalgebra.
\end{theorem}

So an even derivation of a transposed Poisson superalgebra produces a ternary
bracket that satisfies the Filippov--Jacobi identity
\cite[p.~495]{AbramovSovetnikov2025}. The paper also shows that the even
derivations of such an algebra form a Jordan algebra. The odd ones close on the
same bracket exactly when the derivation squares to zero
\cite[Theorem~2, p.~492]{AbramovSovetnikov2025}.

The next section keeps
the Lie product and twists the Jacobi identity by an algebra endomorphism:
Hom-Lie algebras and Lie conformal algebras. % claim-ok: own linking sentence to Section 7

The extensions of the Poisson bracket to differential forms and multivector fields are described~\cite{WlExtensionsOfThePoissonBracketT2020}.
